\documentclass[12pt]{article}
\usepackage{styles/cwilliams-standard}
\usepackage{tablefootnote}
\usepackage{tcolorbox}
\usepackage{multirow}
\usepackage{listings}
\lstset{
    basicstyle=\ttfamily\small,
    numbers=left,
    numberstyle=\tiny\color{gray},
    numbersep=8pt,
    breaklines=true,
    frame=single,
    rulecolor=\color{gray!40},
    backgroundcolor=\color{gray!6},
    keywordstyle=\color{blue!70},
    stringstyle=\color{orange!80!black},
    showstringspaces=false,
    language=Python,
}

\setclass{\NLP}
\settitle{Homework 2: Affect, Perspective, and Subective Judgement}

\definecolor{tableheader}{rgb}{0.18, 0.32, 0.55}
\definecolor{rowalt}{rgb}{0.93, 0.95, 0.98}
\definecolor{modelheader}{rgb}{0.05, 0.42, 0.42}
\definecolor{modelhighlight}{rgb}{1.0, 0.88, 0.50}
\definecolor{llmheader}{rgb}{0.45, 0.08, 0.15}
\definecolor{llmrowalt}{rgb}{0.98, 0.94, 0.94}
\definecolor{llmhighlight}{rgb}{0.68, 0.92, 0.78}

\begin{document}

\maketitlepage

\begin{important}{Provided Code}
    The code files provided are meant to be looked over as Task1.py, Task2.py, and 
    Task3\_Cerebras.py. The other Task3.py is the method I used to query Gemini, but 
    that is not where (the majority of) my examples come from.
\end{important}

\tableofcontents

\section{Preprocessing}

All three tasks share a common data loading step applied to reviewer.csv. The file contains 
a saved row index as an unnamed first column, which is absorbed with index\_col=0 to prevent 
it from being treated as a data column. Review bodies are decoded for HTML entities 
(e.g., \texttt{\&\#39; → '}) using Python's \texttt{html.unescape()}, as the raw CSV contains browser-escaped 
characters from the original web scrape. Rows whose Victim/Abuser/Third label is not one of 
A, V, or T, primarily P (Potential) entries, are dropped, as they do not correspond to 
any of the three target classes.

\section{Affective Analysis}

\begin{enumerate}
    \item \textbf{Leveraging lexicons provided by Mohammed et al.}
    
    To ``leverage'' the lexicons provided, I imported the NRC-VAD-Lexicon-v2.1 from 
    this path: \texttt{NRC-VAD-Lexicon-v2.1/Unigrams/unigrams-NRC-VAD-Lexicon-v2.1.txt}. 
    
    I loaded this into a dictionary where the term is the key, and the valence, arousal, dominance (VAD) are values. 
    This allows for a faster lookup (O(1)) and quicker processing down the line. Given that the 
    downloaded files already have VAD provided, there was no computation that 
    really needed to be done.

    \item \textbf{Relevant Sentence Extraction}
    
    Given the three focus words (spy, stalk, stealth), I went through all sentences in the dataset
    \texttt{reviewer.csv} and used the spaCy \texttt{en\_core\_web\_sm} package to seperate each word in each of the sentences. 
    Looping through the split sentences, I then check if the sentence contains one of the three focus words. 
    If it does contain the focus word, save the entire sentence along with the reviewer type, add that to a list 
    as a tuple, and the move on. 

    \begin{table}[H]
        \centering
        \renewcommand{\arraystretch}{1.4}
        \begin{tabular}{p{10cm} p{3cm}}
            \toprule
            \rowcolor{tableheader}
            \color{white}\textbf{Sentence} & \color{white}\textbf{Reviewer} \\
            \midrule
            \rowcolor{rowalt} I love that you can take a picture and am constantly spying on coworkers at my desk from my laptop. & abuser \\
            This is also good to spy on someone. & third\_party \\
            \rowcolor{rowalt} i can spy on my child whenever i want its amazing he cant go anywhere without me knowing look & abuser \\
            Great app for families keeping track of their kids, friends or if they just wanna be creepy stalkers. & abuser \\
            \rowcolor{rowalt} Another way of spying on people. & abuser \\
            \bottomrule
        \end{tabular}
        \caption{Example Sentences Containing Focus Words}
        \label{tab:sentences}
    \end{table}

    \item \textbf{Adjective Extraction}
    
    For this subtask, I went through all sentences in the \texttt{reviewer.csv} and used the spaCy 
    \texttt{en\_core\_web\_sm} package again to POS tag every word. Once every word is tagged, go through each 
    sentence and then each word, and if the word's POS tag is ``\texttt{ADJ}'', then add that word and it's 
    corresponding VAD scores to a dictionary where the key is the reviewer type, and the 
    values are a list of dictionary's that have the ADJ's as keys and the values are their VAD scores.

\begin{lstlisting}
{   'third_party': [
        {'word': 'good', 
        'valence': 0.876, 'arousal': -0.264, 'dominance': 0.068},
        ...
    ],
    'abuser': [
        {'word': 'amazing',
         'valence': 0.812, 'arousal': 0.674, 'dominance': 0.698}, 
        ...
    ],
    'victim': [
        {'word': 'teen',
         'valence': 0.458, 'arousal': 0.062, 'dominance': 0.136}, 
         ...
    ]
}
\end{lstlisting}

    \item \textbf{Average scores of adjectives for all dimensions}
    
    Using the dictionary created in the previous subtask, I go by key (reviewer type) and get all scores and average them
    across the VAD domains. This is rather straight forward, and here are the scores that I got for this task.

    \begin{table}[H]
        \centering
        \renewcommand{\arraystretch}{1.4}
        \begin{tabular}{p{4cm} p{3cm} p{3cm} p{3cm}}
            \toprule
            \rowcolor{tableheader}
            \color{white}\textbf{Reviewer} & \color{white}\textbf{Valence} & \color{white}\textbf{Arousal} & \color{white}\textbf{Dominance} \\
            \midrule
            \rowcolor{rowalt} third\_party & 0.186 & -0.006 & 0.055 \\
                              abuser & 0.498 & -0.005 & 0.237 \\
            \rowcolor{rowalt} victim & 0.100 & 0.055 & 0.034 \\
            \bottomrule
        \end{tabular}
        \caption{Average VAD Scores by Reviewer Type}
        \label{tab:vad-averages}
    \end{table}

    \item \textbf{Compare dimensions}
    
    As per the table in the previous subtask, we can see there are values per reviewer type that 
    describe emotions and intensity of emotions. 

    \begin{table}[H]
        \centering
        \renewcommand{\arraystretch}{1.4}
        \begin{tabular}{p{4cm} p{3cm} p{3cm} p{3cm}}
            \toprule
            \rowcolor{tableheader}
            \color{white}\textbf{Reviewer} & \color{white}\textbf{Valence} & \color{white}\textbf{Arousal} & \color{white}\textbf{Dominance} \\
            \midrule
            \rowcolor{rowalt} third\_party & 0.186 & \cellcolor{red!13}\textbf{-0.006} & 0.055 \\
                              abuser & \cellcolor{green!13}\textbf{0.498} & -0.005 & \cellcolor{green!13}\textbf{0.237} \\
            \rowcolor{rowalt} victim & \cellcolor{red!13}\textbf{0.100} & \cellcolor{green!13}\textbf{0.055} & \cellcolor{red!13}\textbf{0.034} \\
            \bottomrule
        \end{tabular}
        \caption{Average VAD Scores by Reviewer Type w/ Highest Value Highlighted. 
        The scores highlighted red are the lowest for that VAD domain. 
        The scores highlighted green are the highest values for that VAD domain}
        \label{tab:vad-averages-highlighted}
    \end{table}

    Abusers appear to have very positive feelings towards ($0.498$ valence) the apps they review and feel a higher 
    degree of control with a dominance of $0.237$ on average when reviewing about the app, showing that they feel they 
    have control over who they are ``stalking'' or abusing. Conversely, victims have less positive feelings (although not negative) 
    about the apps but feel more intensely about the apps with an arousal of $0.055$ on average and feel much less in control than 
    the abusers with a dominance of $0.034$. These two go hand in hand and it makes sense overall that the third\_party reviwers do not have particular strong 
    feelings towards this app as they are not the ones using the apps. 

    Third\_party reviewers have more positive feelings about this than victims which does make sense given that 
    they are not experiencing the effects of these apps and are only commenting on them. And, they seem to feel 
    more in control than victims which also makes sense.

    Ultimately, it appears that abusers feel the most in control in relation to these apps, victims feel the 
    least positive and in control, and third\_party is in between.

\end{enumerate}

\section{Joint Modeling for Severity and Convincingness}

For this task, I selected BERT and DistilBERT for the two models that I will be comparing. 

For Task 2, the continuous severity and convincingness scores (1–5 scale) are discretized 
into three ordinal bins: low ($\leq 2$), medium ($2 < s \leq 3$), and high ($> 3$). Each 
review is then encoded as a 6-dimensional binary label vector combining both tasks, and 
tokenized using each model's pretrained tokenizer with truncation and padding to 128 tokens.

\begin{table}[H]
    \centering
    \renewcommand{\arraystretch}{1.4}
    \begin{tabular}{l l c c c c}
        \toprule
        \rowcolor{modelheader}
        \color{white}\textbf{Model} & \color{white}\textbf{Type} & \color{white}\textbf{Precision} & \color{white}\textbf{Recall} & \color{white}\textbf{F1} & \color{white}\textbf{Macro F1} \\
        \midrule
        \multirow{2}{*}{BERT}       & Severity       & \cellcolor{modelhighlight}\textbf{0.7480} & \cellcolor{modelhighlight}\textbf{0.7401} & \cellcolor{modelhighlight}\textbf{0.7336} & \multirow{2}{*}{0.6755} \\
                                     & Convincingness & 0.6392 & 0.6378 & 0.6174 & \\
        \cmidrule{1-6}
                          \multirow{2}{*}{DistilBERT} & Severity       & 0.7140 & 0.7280 & 0.7190 & \multirow{2}{*}{0.6817} \\
                                                       & Convincingness & \cellcolor{modelhighlight}\textbf{0.6978} & \cellcolor{modelhighlight}\textbf{0.6543} & \cellcolor{modelhighlight}\textbf{0.6444} & \\
        \bottomrule
    \end{tabular}
    \caption{BERT vs.\ DistilBERT Classification Performance. Amber highlighted cells show the best
    performance for that type across models. One cell could be read as ``BERT has the best score for
    precision in severity across both models''}
    \label{tab:model-comparison}
\end{table}

Table~\ref{tab:model-comparison} reports Precision, Recall, and F1 for each task and Macro F1 across
5-fold cross-validation. Both models predict severity more reliably than convincingness, 
BERT achieves F1 of 0.7336 vs.\ 0.6174, and DistilBERT 0.7190 vs.\ 0.6444. This gap likely reflects 
that severity is more directly signaled by lexical features (explicit threat or harm language), 
whereas convincingness is a more subjective property requiring inference beyond surface wording.

BERT outperforms DistilBERT on severity across all three metrics (P: 0.7480 vs.\ 0.7140, R: 0.7401 
vs.\ 0.7280, F1: 0.7336 vs.\ 0.7190), suggesting its larger capacity better captures the linguistic 
markers that distinguish severity levels. However, DistilBERT outperforms BERT on convincingness 
(F1: 0.6444 vs.\ 0.6174) and achieves a higher overall Macro F1 (0.6817 vs.\ 0.6755). The task gap 
is also narrower for DistilBERT (0.0746) than for BERT (0.1162), indicating more balanced joint learning.

I prefer \textbf{DistilBERT} for this task. Although BERT achieves the best individual score on severity, the 
task of joint prediction, where both targets must be predicted well simultaneously, is more suited to 
DistilBERT. The models higher Macro F1 and smaller inter-task performance gap make it the stronger 
joint model. Also, DistilBERT has 40\% fewer parameters than BERT, trainig approximately twice as 
fast with only marginal accuracy loss, which is an advantage (in my book) in a cross-validation schema.

\section{Role-Aware LLM Analysis}

\begin{important}
    For all prompts, Cerebras and Gemini, I used \textit{Chain of Thought} prompting.
\end{important}

For Task 3, the labeled portion of reviewer.csv is split into a 203-sample pool for 
few-shot example selection and a 200-sample held-out test set. For cosine similarity 
retrieval, all pool and test texts are jointly fit to a TF-IDF vectorizer with a 
5,000-feature vocabulary and English stop word removal, producing sparse document 
vectors used to compute pairwise similarity at inference time.

I initially used the assignment given Gemini 2.5 Flash Lite model. Given some 
restrictions on the rate limits, I found that using the Gemini 3 Flash (preview) model allowed for 
more requests, and the TA (thank you Nikhil) recommend using Cerebras Llama 3.1-8B which had unlimited 
(or for our use case \textit{essentially} unlimited) calls.

When I began API calls nearly a week and some in advance, the rate limits were going to make this difficult. The RPD 
(requests per day) for the Gemini 2.5 Flash Lite was said to be 1000. A potential work around I thought of was to use 
three API keys (on separate accounts) to grant myself higher rates because I was getting many 429 errors (too busy, too many 
people using the API). But, it turned out, that 2.5 has a rate limit of \textbf{20} requests per day. That was not 
going to cut it. So, I switched to Gemini 3 Flash (3 API Keys), and then Cerebras (1 API Key).

Long story short, this section will have the results of the Cerebras model with some minor 
comparisons between Gemini 3 and Llama 3.1 to see if there is a strong difference between the 
model sizes in the role aware analysis.

\subsection{Zero Shot \& Few-Shot Detection (Gemini and Llama)}

Here is the system prompt that I used:

\begin{tcolorbox}[
    title=System Prompt,
    colback=llmrowalt,
    colframe=llmheader,
    colbacktitle=llmheader,
    coltitle=white,
    fonttitle=\bfseries\small,
    fontupper=\ttfamily\small,
    boxrule=0.6pt,
    arc=3pt,
]
You are a text classifier. Given an app review, classify the narrator's role as one of:

\smallskip
\hspace*{1em}\textbf{Abuser (A):} The reviewer is using or endorses using the app to spy on, stalk, or control someone.\\
\hspace*{1em}\textbf{Victim (V):} The reviewer is someone being spied on, stalked, or controlled through the app.\\
\hspace*{1em}\textbf{Third-party (T):} The reviewer is a neutral observer commenting on the app without being an abuser or victim.

\smallskip
Think step by step about who the narrator is and their relationship to surveillance/stalking.

After your reasoning, write your final answer on the last line as: 

\textbf{ANSWER: A}, \textbf{ANSWER: V}, or \textbf{ANSWER: T}.
\end{tcolorbox}

Every system prompt was then either appended with the example that it needs to classify (zero-shot)
or a few examples and then the example it needs to classify (few-shot)

\vspace{0.5em}
\begin{tcbraster}[
    raster columns=2,
    raster column skip=0.04\textwidth,
    raster valign=top,
    colback=llmrowalt,
    colframe=llmheader,
    colbacktitle=llmheader,
    coltitle=white,
    fonttitle=\bfseries\small,
    fontupper=\ttfamily\small,
    boxrule=0.6pt,
    arc=3pt,
]
\begin{tcolorbox}[title=Zero-Shot]
Classify this review: Awesome app can spy on people or set up in my front window
\end{tcolorbox}
\begin{tcolorbox}[title=Few-Shot]
Here are some examples:

\textbf{Example 1:}

\textbf{Review:} Awesome app can spy on people or set up in my front window

\textbf{Label:} A

\textbf{Example 2:}

\textbf{Review:} Now spying on my roommate is one of my 79 daily tasks.

\textbf{Label:} A

Now classify this review: ``This is the worst update ever! Why would you do this?!
It was the best app, now I can't stalk anymore lol''
\end{tcolorbox}
\end{tcbraster}

\vspace{0.5em}

With these two prompts, this got all my results. This was the format for all queries to 
Llama 3.1 - 8B and Gemini 3 Flash. The Few-Shot prompting was cycled through 
with multiple $k$ values (the number of examples) (the example above is K = 2) 
for all hand picked, random, and cosine chosen examples. For cosine few-shot specifically, 
I created TF-IDF vectors for all reviews, and then find the K closest TF-IDF vectors of 
other reviews to the test example (i.e., find the 2 closest reviews to the review we 
want to classify).  

\vspace{0.5em}

\begin{table}[H]
    \centering
    \renewcommand{\arraystretch}{1.4}
    \begin{tabular}{l l c c c c}
        \toprule
        \rowcolor{llmheader}
        \color{white}\textbf{Experiment} & \color{white}\textbf{Class} & \color{white}\textbf{Precision} & \color{white}\textbf{Recall} & \color{white}\textbf{F1} & \color{white}\textbf{Macro F1} \\
        \midrule
        \multirow{3}{*}{Zero-shot}
            & Abuser      & 0.8077 & 0.8750 & 0.8400 & \multirow{3}{*}{0.5014} \\
            & Victim      & 0.5179 & 0.6744 & 0.5859 & \\
            & Third-party & 0.1429 & 0.0541 & 0.0784 & \\
        \cmidrule{1-6}
        \multirow{3}{*}{Few-shot Handpicked ($k$=2)}
            & Abuser      & 0.8417 & 0.8417 & 0.8417 & \multirow{3}{*}{0.5786} \\
            & Victim      & \cellcolor{llmhighlight}\textbf{0.5652} & 0.6047 & 0.5843 & \\
            & Third-party & \cellcolor{llmhighlight}\textbf{0.3235} & \cellcolor{llmhighlight}\textbf{0.2973} & \cellcolor{llmhighlight}\textbf{0.3099} & \\
        \cmidrule{1-6}
        \multirow{3}{*}{Few-shot Random ($k$=5)}
            & Abuser      & 0.8092 & \cellcolor{llmhighlight}\textbf{0.8833} & 0.8446 & \multirow{3}{*}{0.5465} \\
            & Victim      & 0.5400 & 0.6279 & 0.5806 & \\
            & Third-party & 0.3158 & 0.1622 & 0.2143 & \\
        \cmidrule{1-6}
        \multirow{3}{*}{Few-shot Cosine ($k$=5)}
            & Abuser      & \cellcolor{llmhighlight}\textbf{0.8750} & 0.8167 & \cellcolor{llmhighlight}\textbf{0.8448} & \multirow{3}{*}{0.5733} \\
            & Victim      & 0.5362 & \cellcolor{llmhighlight}\textbf{0.8605} & \cellcolor{llmhighlight}\textbf{0.6607} & \\
            & Third-party & 0.3158 & 0.1622 & 0.2143 & \\
        \bottomrule
    \end{tabular}
    \caption{Per-class performance for key Cerebras (Llama 3.1-8B) experiments.
    Mint cells show the best value for that class and metric across all four experiments. e.g., 
    the best recall for Abuser is the Few-shot Random ($k$=5).}
    \label{tab:Cerebras-perclass}
\end{table}


Table~\ref{tab:Cerebras-perclass} reveals a pattern driven primarily by class imbalance: 
with Abuser comprising 60\% of the test set, all four experiments classify it reliably (F1 0.840--0.845), 
while the minority classes Victim (21.5\%) and Third-party (18.5\%) show weaker classification.

Zero-shot performs competitively on Abuser but nearly fails on Third-party (F1=0.078, recall=0.054). 
Without examples to anchor minority class boundaries, the model defaults heavily toward predicting
 Abuser, treating ambiguous reviews as the majority class by default.

Few-shot Handpicked ($k$=2) produces the most balanced results. Introducing just two curated examples 
raises Third-party F1 from 0.078 to 0.310 — a fourfold improvement and the largest single gain 
across all cells in the table. This is because Third-party is a subtle category requiring a concrete 
demonstration of neutral observation, which curated examples directly provide. Handpicked also achieves 
the highest Victim precision (0.5652) and the best Macro F1 overall (0.5786), making it the strongest 
general-purpose strategy.

Few-shot Random ($k$=5) achieves the highest Abuser recall (0.8833), slightly outperforming 
zero-shot's 0.8750. Stratified random sampling ensures exposure to a variety of Abuser phrasings, 
which sharpens the model's sensitivity to that class. However, Random offers no improvement over Cosine 
on Third-party — both score identically (F1=0.214) — suggesting that at $k$=5, randomly selected 
examples are no less informative than similarity-retrieved ones for that class. This likely 
reflects the small Third-party pool size: cosine retrieval frequently returns Abuser or Victim 
neighbors regardless, providing no class-specific signal beyond what random selection would offer.

Few-shot Cosine ($k$=5) takes a different trade-off. It achieves the best Victim recall 
(0.8605) and F1 (0.6607) by retrieving the most lexically similar pool examples for each test 
review. Victim reviews are contextually diverse — a parent tracking a child, a partner being 
monitored, an employee under surveillance — so finding a similar labeled example per test 
case anchors the prediction more effectively than fixed handpicked or random ones. However, 
this benefit does not extend to Third-party, where Handpicked remains superior.

Overall, no single strategy dominates across all classes. Handpicked ($k$=2) is preferred when 
balanced class coverage and Macro F1 are the priority, while Cosine ($k$=5) is the better choice 
when Victim recall is critical.

\subsection{Comparing Performance}

\begin{table}[H]
    \centering
    \renewcommand{\arraystretch}{1.4}
    \begin{tabular}{l l c c c c}
        \toprule
        \rowcolor{llmheader}
        \color{white}\textbf{Strategy} & \color{white}\textbf{Class} & \color{white}\textbf{Precision} & \color{white}\textbf{Recall} & \color{white}\textbf{F1} & \color{white}\textbf{Macro F1} \\
        \midrule
        \multirow{3}{*}{Zero-shot}
            & Abuser      & 0.8077 & 0.8750 & 0.8400 & \multirow{3}{*}{0.5014} \\
            & Victim      & 0.5179 & 0.6744 & 0.5859 & \\
            & Third-party & 0.1429 & 0.0541 & 0.0784 & \\
        \cmidrule{1-6}
        \multirow{3}{*}{Handpicked ($k$=2)}
            & Abuser      & 0.8417 & 0.8417 & 0.8417 & \multirow{3}{*}{\textbf{0.5786}} \\
            & Victim      & \cellcolor{llmhighlight}\textbf{0.5652} & 0.6047 & 0.5843 & \\
            & Third-party & \cellcolor{llmhighlight}\textbf{0.3235} & \cellcolor{llmhighlight}\textbf{0.2973} & \cellcolor{llmhighlight}\textbf{0.3099} & \\
        \cmidrule{1-6}
        \multirow{3}{*}{Random ($k$=5)}
            & Abuser      & 0.8092 & \cellcolor{llmhighlight}\textbf{0.8833} & 0.8446 & \multirow{3}{*}{0.5465} \\
            & Victim      & 0.5400 & 0.6279 & 0.5806 & \\
            & Third-party & 0.3158 & 0.1622 & 0.2143 & \\
        \cmidrule{1-6}
        \multirow{3}{*}{Cosine ($k$=5)}
            & Abuser      & \cellcolor{llmhighlight}\textbf{0.8750} & 0.8167 & \cellcolor{llmhighlight}\textbf{0.8448} & \multirow{3}{*}{0.5733} \\
            & Victim      & 0.5362 & \cellcolor{llmhighlight}\textbf{0.8605} & \cellcolor{llmhighlight}\textbf{0.6607} & \\
            & Third-party & 0.3158 & 0.1622 & 0.2143 & \\
        \bottomrule
    \end{tabular}
    \caption{Cerebras (Llama 3.1-8B) per-class performance by example selection strategy
    at best $k$. Mint cells show the best value for that class and metric across all strategies.}
    \label{tab:Cerebras-strategy}
\end{table}


We can see that the Zero-shot method did not win in any of the strategies, and this stands to reason 
that without examples, the LLM would have a lot of difficultly learning its task given its prompt. 
It's also important to point out (for myself) that the Handpicked examples achieved the highest 
Macro-F1 and greatest amount of classification for Third-party and precision for victim. It is 
also impressive that handpicked peaked at $k$=2 whereas Random and Cosine had to go to $k$=5 suggests that 
quality matters more than quantity. Two carefully chosen examples outperform five randomly or 
automatically retrieved ones in overall accuracy. 

Random sits between the two in F1 (0.5465) but requires $k$=5 to reach that ceiling, 
making it the least efficient strategy. The narrow gap between Handpicked (0.5786) and Cosine 
(0.5733) at their respective best $k$ values indicates that both approaches are competitive 
at the macro level, with the distinction emerging only in per-class analysis where their 
trade-offs become apparent.

\begin{table}[H]
    \centering
    \renewcommand{\arraystretch}{1.4}
    \begin{tabular}{l l c c c c}
        \toprule
        \rowcolor{llmheader}
        \color{white}\textbf{Model} & \color{white}\textbf{Class} & \color{white}\textbf{Precision} & \color{white}\textbf{Recall} & \color{white}\textbf{F1} & \color{white}\textbf{Macro F1} \\
        \midrule
        \multirow{3}{*}{Gemini 3 Flash (Preview)}
            & Abuser      & \cellcolor{llmhighlight}\textbf{0.8250} & 0.8250 & 0.8250 &  \\
            & Victim      & \cellcolor{llmhighlight}\textbf{0.5800} & \cellcolor{llmhighlight}\textbf{0.6744} & \cellcolor{llmhighlight}\textbf{0.6237} & \textbf{0.5923} \\
            & Third-party & \cellcolor{llmhighlight}\textbf{0.3667} & \cellcolor{llmhighlight}\textbf{0.2973} & \cellcolor{llmhighlight}\textbf{0.3284} & \\
        \cmidrule{1-6}
        \multirow{3}{*}{Cerebras Llama 3.1-8B}
            & Abuser      & 0.8077 & \cellcolor{llmhighlight}\textbf{0.8750} & \cellcolor{llmhighlight}\textbf{0.8400} & \multirow{3}{*}{0.5014} \\
            & Victim      & 0.5179 & \cellcolor{llmhighlight}\textbf{0.6744} & 0.5859 & \\
            & Third-party & 0.1429 & 0.0541 & 0.0784 & \\
        \bottomrule
    \end{tabular}
    \caption{Zero-shot performance comparison between Gemini 3 Flash (Preview) and Cerebras
    Llama 3.1-8B. Mint cells show the best value per class and metric across both models.
    V Recall is tied at 0.6744.}
    \label{tab:model-zeroshot}
\end{table}

This is the most impressive table from my experiments with prompting the LLM's. Given that rate limits 
were making API calling slow and difficult, I was only able to complete Zero-shot calling Gemini. 
The Few-shot experiments were incomplete (even with 3 API keys) and did not finish, so I can not reasonably 
make comparisons between the Few-shot experiments.

Within zero-shot, Gemini 3 Flash outperforms Cerebras Llama 3.1-8B overall 
(Macro F1: 0.5923 vs.\ 0.5014), with the gap driven almost entirely by Third-party detection. Gemini 
achieves a Third-party F1 of 0.3284 compared to Cerebras's 0.0784, a fourfold difference, suggesting 
that the larger, more instruction-tuned model better recognizes neutral observers without any examples 
to guide it. This is consistent with the expectation that stronger zero-shot instruction-following 
correlates with model scale and alignment training.

Cerebras, however, outperforms Gemini on Abuser Recall (0.8750 vs.\ 0.8250) and Abuser F1 
(0.8400 vs.\ 0.8250). This likely reflects the tendency of smaller models to default toward the 
majority class, which is a bias that incidentally helps recall on Abuser but collapses performance on 
Third-party. Both models tie exactly on Victim Recall (0.6744), suggesting similar sensitivity to 
victim-indicating language even across a large difference in model size.

\section{Ethics}

This work involves sensitive subject matter: app reviews written in the context of 
surveillance, stalking, and domestic control. This raises several ethical 
considerations that informed both the methodology and interpretation of results.

The NRC VAD Lexicon used in Task 1 was itself constructed with ethical safeguards: 
approximately 200 entries associated with identity group terms, slurs, or taboo 
language were removed by the lexicon authors to prevent harmful associations 
from propagating into downstream analyses. As noted in the lexicon's documentation, 
affective scores reflect historical annotator perceptions and may embed cultural 
biases, particularly for terms associated with marginalized groups. Results from 
Task 1 should therefore be interpreted as reflecting patterns in reviewer 
language rather than objective affective ground truth.

In Task 3, reviews are sent to third-party API providers (Google Gemini 
and Cerebras) for classification. This means potentially sensitive review
content including descriptions of stalking, spousal surveillance, and 
child tracking is transmitted to external servers. In a production or 
clinical setting, this would require explicit data handling agreements 
and user consent.

The classification task itself carries asymmetric risks of misclassification. 
Labeling a victim as an abuser, or an abuser as a third party, could have 
serious consequences in any real-world application of such a system. The 
models consistently underperform on the Victim and Third-party classes 
due to label imbalance, meaning the groups most in need of accurate 
identification are precisely those the models handle least reliably. 
Any deployment of these classifiers would require either rebalancing 
strategies or human review of low-confidence predictions before any 
consequential action is taken.

\section{Code}

\begin{aside}{Provided Code}
    The code is provided with this submission on Canvas, but in case something is not correct 
    or does not work, \href{https://drive.google.com/drive/folders/1W0iUTP5_-hPrcNcUi155nI08rTYZorFQ?usp=sharing}{here} is the link to the files on Google Drive.
\end{aside}

\end{document}